\subsection{EXPONENTIATION OF CARDINALS}

DEFINITION 4. Let a and b be cardinals. The cardinal of the set of mappings of b into a is denoted by \(\mathfrak{a}^{\mathfrak{b}}\) , by abuse of notation.

The abuse of notation here lies in the fact that \(a^{b}\) already denotes the set of graphs of mappings of b into a (Chapter II, §5, no. 3), and this set is not necessarily a cardinal (Exercise 2). It will always be clear from the context which meaning is to be attached to the symbol \(a^{b}\) .

PROPOSITION 9. Let X and Y be two sets, a and b their respective cardinals. Then the set \(X^{Y}\) has cardinal \(a^{b}\) .

For there exists a bijection of \(\mathbf{X}^{\mathrm{Y}}\) onto the set of mappings of \(\mathfrak{b}\) into \(\mathfrak{a}\) (Chapter II, § 5, no. 2, Proposition 2, Corollary).

PROPOSITION 10. Let \(\mathfrak{a}\) and \(\mathfrak{b}\) be cardinals and let \(\mathbf{I}\) be a set such that Card (I) = \(\mathfrak{b}\) . If \(\mathfrak{a}_{\iota} = \mathfrak{a}\) for all \(\iota \in \mathbf{I}\) , we have \(\mathfrak{a}^{\mathfrak{b}} = \underset{\iota \in \mathbf{I}}{\mathsf{P}}\mathfrak{a}_{\iota}\) .

This follows from the definition of the product of a family of sets as a set of functional graphs (Chapter II, § 5, no. 3).

COROLLARY 1. Let \(\mathfrak{a}\) be a cardinal and let \((\mathfrak{b}_t)_{t\in I}\) be a family of cardinals. Then

\[
\mathfrak{a}^{\sum_{i \in I} \mathfrak{b}_i} = \underset {i \in I} {\mathbf {P}} \mathfrak{a}^{\mathfrak{b}_i}.
\]

Let S be the sum of the sets \(\mathfrak{b}_i\) , and put \(a_s = a\) for all \(s \in S\) . Both sides of the equality to be proved are then equal to \(\underset{s \in S}{\mathsf{P}} a_s\) , by virtue of Proposition 10 and the associativity formula for products (no. 3, Proposition 5(b)).

COROLLARY 2. Let \((\mathfrak{a}_i)_{i\in I}\) be a family of cardinals and let \(\mathfrak{b}\) be a cardinal. Then

\[
\left(\underset {i \in I} {\mathbf {P}} a _ {i}\right) ^ {b} = \underset {i \in I} {\mathbf {P}} a _ {i} ^ {b}.
\]

Let \(a_{i\beta} = a_i\) for each pair \((i, \beta) \in I \times b\) . Then, by associativity of the product, we have

\[
\left(\underset {i \in I} {\mathsf {P}} a _ {i}\right) ^ {\mathfrak {b}} = \underset {\beta \in \mathfrak {b}} {\mathsf {P}} \left(\underset {i \in I} {\mathsf {P}} a _ {i \beta}\right) = \underset {i \in I} {\mathsf {P}} \left(\underset {\beta \in \mathfrak {b}} {\mathsf {P}} a _ {i \beta}\right) = \underset {i \in I} {\mathsf {P}} a _ {i} ^ {\mathfrak {b}}.
\]

COROLLARY 3. Let \(a, b, c\) be cardinals. Then \(a^{bc} = (a^{b})^{c}\) .

Let \(\mathfrak{b}_{\gamma} = \mathfrak{b}\) for all \(\gamma \in \mathfrak{c}\) . Then

\[
a ^ {b c} = a ^ {\sum_ {\gamma \in c} b _ {\gamma}} = \underset {\gamma \in c} {\mathbf {P}} a ^ {b _ {\gamma}} = (a ^ {b}) ^ {c}
\]

by virtue of Corollary 1.

PROPOSITION 11. Let \(\mathfrak{a}\) be a cardinal. Then \(\mathfrak{a}^0 = 1\) , \(\mathfrak{a}^1 = \mathfrak{a}\) , \(1^{\mathfrak{a}} = 1\) ; and \(0^{\mathfrak{a}} = 0\) if \(\mathfrak{a} \neq 0\) .

For there exists a unique mapping of \(\varnothing\) into any given set (namely, the mapping whose graph is the empty set); the set of mappings of a set consisting of a single element into an arbitrary set X is equipotent to X (Chapter II, §5, no. 3); there exists a unique mapping of an arbitrary set into a set consisting of a single element; and, finally, there is no mapping of a non-empty set into \(\varnothing\) .

Note in particular that \(0^{0}=1\) .

PROPOSITION 12. Let X be a set and let a be its cardinal. Then the cardinal of the set \(\mathfrak{P}(X)\) of all subsets of \(X\) is \(2^{\mathfrak{a}}\) .

Let \(\alpha\) and \(\beta\) be the elements of the cardinal 2. For each subset Y of X let \(f_{\mathbf{Y}}\) be the mapping of X into 2 defined by \(f_{\mathbf{Y}}(x) = \alpha\) if \(x \in \mathbf{Y}\) and \(f_{\mathbf{Y}}(x) = \beta\) if \(x \in \mathbf{X} - \mathbf{Y}\) . Let \(u\) be the mapping \(\mathbf{Y} \to f_{\mathbf{Y}}\) of \(\mathfrak{P}(\mathbf{X})\) into \(2^{\mathbf{x}}\) . Conversely, with each mapping \(g\) of X into 2 let us associate the subset \(-g^{1}(\alpha)\) of X, and let \(v\) be the mapping \(g \to -g^{1}(\alpha)\) of \(2^{\mathbf{x}}\) into \(\mathfrak{P}(\mathbf{X})\) . It is clear that \(u \circ v\) and \(v \circ u\) are the identity mappings of \(2^{\mathbf{x}}\) and \(\mathfrak{P}(\mathbf{X})\) ; hence \(u\) and \(v\) are bijections (Chapter II, §3, no. 8, Proposition 8, Corollary) and therefore Card \((\mathfrak{P}(\mathbf{X})) = 2^{\mathbf{a}}\) .

